\section{Lower bounds for relational circuits}\label{sec:logical}
The construction of \S\ref{sec:preservation} is independent of the
query. We now interpret the input relations of $G_k$ on it; the read
edge of each test identifies the witness color that the fold removes.

\subsection{The tests for \texorpdfstring{$G_k$}{Gk}}\label{sec:G}
On a color $\mu$, interpret $P_i$ as $1$ precisely when $i\in\mu$, and
$Q_i$ as $1$ precisely when $i\notin\mu$. Put both values zero on every
loop. More explicitly,
\[
 P_i^{\mathcal M_0}\coloneqq\bigcup_{\mu\in\Lambda:\,i\in\mu}e_\mu,\qquad
 Q_i^{\mathcal M_0}\coloneqq\bigcup_{\mu\in\Lambda:\,i\notin\mu}e_\mu.
\]
For each color $\lambda$, use the read edge $(\hub,v_{\lambda^c})$,
of color $\lambda^c$ (\S\ref{sec:structure}), and the fold of $\lambda$
(\S\ref{sec:types}). Since $\lambda\ne\lambda^c$, the fold does not
recolor this read edge.

\begin{lemma}[The tests for $G_k$]\label{lem:G-tests}
For every color $\lambda$,
\[
 \sem{G_k}{\mathcal M_0}(\hub,v_{\lambda^c})=1,\qquad
 \sem{G_k}{\mathcal M_\lambda}(\hub,v_{\lambda^c})=0.
\]
Moreover, the sentence $G_k^\forall=\forall x\,\forall y\,(x=y\lor
G_k(x,y))$ of \S\ref{sec:main-results} is true in $\mathcal M_0$
and false in every $\mathcal M_\lambda$. Thus $G_k$ and $G_k^\forall$
each supply $2^k$ tests.
\end{lemma}
\begin{proof}
For non-loop edges with colors $\mu,\nu$, the body of $G_k$ holds
exactly when $\nu=\mu^c$. A loop, whose $P_i,Q_i$ values are all zero,
cannot be a witness edge. Hence, on $\mathcal M_0$ and on every
$\mathcal M_\lambda$, $G_k$ holds at a pair $(u,v)$ of distinct vertices
exactly when $u$ has an edge whose color is the complement of the color
of $(u,v)$. Against the read color $\lambda^c$, the sole accepted
witness color is $\lambda$. That color occurs at $\hub$ in
$\mathcal M_0$ and does not occur there after the fold. In
$\mathcal M_0$, every vertex has edges of every color, so $G_k$ holds at
every pair of distinct vertices, and $x=y$ handles the diagonal; in
$\mathcal M_\lambda$, $G_k$ fails at the pair $(\hub,v_{\lambda^c})$ of
distinct vertices.
\end{proof}

\subsection{Consequences}\label{sec:query-bounds}
\mainbounds*
\begin{proof}
Let $C\eqfin G_k$, and take the reference structure at
any size $n\ge n_0(2^k)$. For each color $\lambda$, the answers of $G_k$
at $(\hub,v_{\lambda^c})$ on $\mathcal M_0$ and on $\mathcal M_\lambda$ differ
by Lemma~\ref{lem:G-tests}, and the fold of $\lambda$ does not recolor
this pair. By Lemma~\ref{lem:unnoticed}, some composition gate of $C$
notices the fold of $\lambda$. By Lemma~\ref{lem:notice}, each
composition gate notices at most two folds, so $2\gcost C\ge2^k$. A term
is a circuit without sharing, so the same bound holds for its number of
compositions.
\end{proof}
Together with the expansion of $G_k$ in \S\ref{sec:main-results}, these
lower bounds show that the minimum circuit size for $G_k$ is
$\Theta(2^k)$; Theorem~\ref{thm:tree} in \S\ref{sec:trees}
determines the minimum term size.

The lower bound does not rest on the difficulty of evaluating $G_k$. As
its proof shows, it holds even for circuits that are correct only on the
reference structure and its folded structures, on which every non-loop
edge has a color from the palette. A read edge there accepts exactly the
witness edges of the complementary color, so a trie of the colors
present in each row evaluates $G_k$ on all pairs of these structures in
time $O(kn^2)$, linear in the size of the input.

\paragraph{Equations.}
By the semantics of \CoR\ equations in \S\ref{sec:prelim-cor},
the same tests also bound the equational form of Tarski's correspondence.
\begin{corollary}[\CoR\ equations]\label{cor:equations}
If a \CoR\ equation $T=S$ holds on exactly those finite structures that
satisfy $G_k^\forall$, then every relational circuit with gates computing
$T$ and $S$ has at least $2^{k-1}-1$ composition gates.
\end{corollary}
\begin{proof}
Let $C$ have gates computing $T$ and $S$, and add gates for
\[
 \mathrm{Eq}\coloneqq\overline{\one;\bigl((T\cap\overline S)\cup(\overline T\cap S)\bigr);\one},
\]
which is the full relation when $T=S$ and empty otherwise. This adds two
compositions, both one-sided. By Lemma~\ref{lem:G-tests}, $G_k^\forall$
is true in $\mathcal M_0$ and false in every $\mathcal M_\lambda$, so
$\mathrm{Eq}$ is full on $\mathcal M_0$ and empty on every
$\mathcal M_\lambda$. Both constants are their own predictions, so these
are $2^k$ tests, and Theorem~\ref{thm:test-cost} with
Lemma~\ref{lem:cor-cost} gives $2^k\le2\gcost C+2$.
\end{proof}
The DNF of \S\ref{sec:main-results} gives the equivalent equation
$\id\cup\bigcup_{A\subseteq[k]}[H_A\cap(H_{A^c};\one)]=\one$
with $2^k$ compositions, so the bound is tight up to a factor of two.

\paragraph{Fixed sizes and bounded error.}
The reference structure and its folded structures share one domain,
which may have any size $n\ge n_0(2^k)$, and the randomized assertion of
Theorem~\ref{thm:test-cost} handles bounded error.
\randomizedbounds*
\begin{proof}
At the prescribed size $n$, Lemma~\ref{lem:palette} gives the reference
structure, and Lemma~\ref{lem:G-tests} gives all $2^k$ tests on that domain.
Apply the expected-cost assertion of Theorem~\ref{thm:test-cost} and
$\tcost(C)\le2\gcost C$. Setting $\varepsilon=0$ includes every exact
size-specific circuit.
\end{proof}

\paragraph{Circuits without converse.}
Without converse, a directed version of the reference structure lets
each composition notice at most one fold, and the bound becomes exact.
\begin{restatable}[Without converse]{corollary}{withoutconverse}\label{cor:without-converse}
For every $k\ge1$, every relational circuit without converse that
defines $G_k$ on finite structures has at least $2^k$ composition gates.
This bound is attained by the DNF.
\end{restatable}
Appendix~\ref{app:directed} proves the corollary with a directed
reference coloring. For full \CoR, Proposition~\ref{prop:promise} gives
a circuit with $2^{k-1}$ compositions that is correct on the reference
structure and all its folded structures. Thus these tests attain the
coefficient two of Lemma~\ref{lem:notice} and cannot give more than
$2^{k-1}$. The optimum for $G_k$ on arbitrary inputs remains between
$2^{k-1}$ and $2^k$.

\subsection{Domain size}\label{sec:domain-size}
On arbitrary inputs, the \defn{profile} $\pi(x,z)\in\{0,1\}^{2k}$ of a
pair lists the values of $P_1,\ldots,P_k,Q_1,\ldots,Q_k$ at $(x,z)$; on
the reference structure, the profile of an edge is determined by its color.
By the definition of $G_k$,
the read edge $(x,y)$ accepts the witness edge $(x,z)$ exactly when, for
every $i$, the profile $\pi(x,y)$ has $P_i$ and $\pi(x,z)$ has $Q_i$, or
$\pi(x,y)$ has $Q_i$ and $\pi(x,z)$ has $P_i$. Hence $G_k(x,y)$ depends
only on $\pi(x,y)$ and the set $\Pi_x\coloneqq\{\pi(x,z):z\in V\}$ of profiles in
row $x$, which has at most $|V|$ elements.

\begin{proposition}[Small domains]\label{prop:small-domains}
For all $k,n\ge1$ there is a converse-free relational circuit with $4kn$
composition gates, all of them row projections $X;\one$, and $O(k^2n^2)$
other gates, that computes $G_k$ on every structure with at most $n$
elements.
\end{proposition}
\begin{proof}
A relation is \defn{row-constant} if its value at $(x,y)$ does not
depend on $y$. By~\eqref{eq:prelim-projections}, each row projection $S;\one$
is row-constant and holds in row $x$ exactly when some pair
$(x,z)$ lies in $S$; Boolean
combinations of row-constant relations are row-constant. For
$1\le b\le2k$, let $E_b$ be the input relation giving bit $b$ of the
profile, and put $E_b^1\coloneqq E_b$ and $E_b^0\coloneqq\overline{E_b}$.

In a structure with at most $n$ elements, call a word
$w\in\{0,1\}^\ell$ \defn{live} in row $x$ if some profile in $\Pi_x$
begins with $w$. Live words of one length begin distinct profiles, so
there are at most $n$ of each length, and the empty word is live because
$\pi(x,x)\in\Pi_x$. The circuit stores the live words of length $\ell$ in
$n$ slots: row-constant relations $L_{\ell,r}$ and $W_{\ell,r,b}$, for
$r\le n$ and $b\le\ell$, state that slot $r$ is occupied in row $x$ and
that its word has bit $b$. Initially $L_{0,1}\coloneqq\one$ and
$L_{0,r}\coloneqq\zero$ for $2\le r\le n$; the only stored word
is empty, so there are no word bits. For each slot $r$ and each
$\sigma\in\{0,1\}$, the relation
\[
 S_{\ell,r,\sigma}\coloneqq L_{\ell,r}\cap
 \bigcap_{b=1}^{\ell}\bigl((E_b\cap W_{\ell,r,b})\cup
     (\overline{E_b}\cap\overline{W_{\ell,r,b}})\bigr)\cap E_{\ell+1}^\sigma
\]
holds at $(x,z)$ exactly when slot $r$ is occupied in row $x$ and
$\pi(x,z)$ begins with its word followed by $\sigma$. The row projection
$S_{\ell,r,\sigma};\one$ therefore tells, in every row at once, whether
this extension is live. The at most $n$ live extensions are placed into the slots of length
$\ell+1$ in the fixed order of their pairs $(r,\sigma)$; no order on
the input vertices is used. Index the $2n$ candidate pairs by
$(r_i,\sigma_i)$, for $1\le i\le 2n$, and put
$f_i\coloneqq S_{\ell,r_i,\sigma_i};\one$. Using only Boolean gates,
compute, for $0\le i\le 2n$ and $0\le j\le n$, the row-constant
relation $C_{i,j}$ that holds in a row exactly when $j$ of
$f_1,\ldots,f_i$ hold:
\[
 C_{0,0}\coloneqq\one,\qquad C_{0,j}\coloneqq\zero\ (j>0),\qquad
 C_{i,j}\coloneqq(\overline{f_i}\cap C_{i-1,j})
       \cup(f_i\cap C_{i-1,j-1}),
\]
where $1\le i\le 2n$, $0\le j\le n$, and
$C_{i,-1}\coloneqq\zero$. For $1\le j\le n$, put
\[
 A_{i,j}\coloneqq f_i\cap C_{i-1,j-1}\qquad(1\le i\le 2n).
\]
In each row, $A_{i,j}$ holds exactly if candidate $i$ is the $j$th
live candidate. In particular, at most one $A_{i,j}$ holds for a given
slot $j$. Thus occupancy and every bit of the new slots are computed by
\[
\begin{aligned}
 L_{\ell+1,j}
   &\coloneqq\bigcup_{i=1}^{2n} A_{i,j},\\
 W_{\ell+1,j,b}
   &\coloneqq\bigcup_{i=1}^{2n}
       (A_{i,j}\cap W_{\ell,r_i,b})
       &&(1\le b\le\ell),\\
 W_{\ell+1,j,\ell+1}
   &\coloneqq\bigcup_{\substack{1\le i\le 2n\\\sigma_i=1}} A_{i,j}.
\end{aligned}
\]
Distinct live prefixes have distinct extensions, and every live word
of length $\ell+1$ extends a live prefix. Thus these formulas preserve
the slot invariant. They use only Boolean operations on row-constant
relations, so no further compositions are needed. At each length, the
$2n$ projections producing
the $f_i$ are the only composition gates. The candidate tests use
$O((\ell+1)n)$ Boolean gates, the prefix counts and ranks use
$O(n^2)$, and updating all slots and their $\ell+1$ bits uses
$O((\ell+1)n^2)$. Summing over the $2k$ lengths, together with the
final output gates below, yields $O(k^2n^2)$ non-composition gates.
After length $2k$ the occupied slots
hold exactly the profiles in $\Pi_x$, and
\[
 G_k=\bigcup_{r=1}^{n}\Bigl(L_{2k,r}\cap\bigcap_{i=1}^k
 \bigl((P_i\cap W_{2k,r,k+i})\cup(Q_i\cap W_{2k,r,i})\bigr)\Bigr),
\]
because the intersection over $i$ says that the read edge $(x,y)$
accepts the witness edges whose profile is in slot $r$. The circuit
makes $2n$ row projections
at each of the $2k$ lengths, $4kn$ in all, and uses no converse.
\end{proof}

\domainsize*
\begin{proof}
By Corollary~\ref{cor:randomized} with $\varepsilon=0$, the lower
bound holds at every domain size $n\ge n_0(2^k)$. Since
\[
n_0(2^k)=\lceil12\cdot4^k(k+1)\ln2\rceil+3\le18k4^k
\]
for all $k\ge1$, we may take $C=18$.
\end{proof}

\subsection{Random inputs}\label{sec:random-inputs}
\begin{corollary}[Random inputs]\label{cor:random-inputs}
Fix $k\ge1$ and $n$. Color the edges of the complete graph on $n$
vertices independently and uniformly with the colors in $2^{[k]}$,
determine $P_i,Q_i$ by the colors as in \S\ref{sec:overview}, and
fix a vertex $\hub$. Choose a color $\lambda$ uniformly, a vertex $v$
uniformly among those with $\chi(\hub,v)=\lambda^c$ (any vertex if there
is none), and a fair bit; if the bit is one, apply the fold of
$\lambda$. Then every relational circuit with $g$ composition gates
answers $G_k(\hub,v)$ correctly with probability at most
$\frac12+g/2^k+\delta$,
where $\delta\coloneqq n^34^k\exp(-(n-3)/4^k)$, and for each
$g\le2^{k-1}$ some circuit with $g$ composition gates answers correctly
with probability at least $\frac12+g/2^k-\delta$.
\end{corollary}
\begin{proof}
Let $\mathcal W$ be the event that the coloring has the two-witness
property. The proof of Lemma~\ref{lem:palette}, with $N=2^k$, gives
$\Pr[\lnot\mathcal W]\le\delta$. Fix a coloring in $\mathcal W$. Then
$\hub$ has edges of every color, so $v$ is uniform among the vertices
with $\chi(\hub,v)=\lambda^c$. For each color $\lambda$, let
$\mathrm{err}_{\lambda,0},\mathrm{err}_{\lambda,1}$ indicate an error of
$C$ at $(\hub,v_{\lambda^c})$ on $\mathcal M_0$ and on $\mathcal M_\lambda$.
The error-count inequality in the proof of Theorem~\ref{thm:test-cost},
with the tests of Lemma~\ref{lem:G-tests} and Lemma~\ref{lem:cor-cost},
gives
\[
 \sum_\lambda(\mathrm{err}_{\lambda,0}+\mathrm{err}_{\lambda,1})
 \ge2^k-\tcost(C)\ge2^k-2g
\]
for every choice of the vertices $v_\mu$ of \S\ref{sec:structure}, hence
also in expectation when each $v_\mu$ is uniform among the vertices with
$\chi(\hub,v_\mu)=\mu$. Given the coloring, the probability of
an error is the average of $\mathrm{err}_{\lambda,b}$ over uniform
$\lambda$ and $b$, so it is
at least $(2^k-2g)/2^{k+1}=\frac12-g/2^k$. Averaging over the colorings
proves the upper bound.

For the lower bound, choose $g$ of the pairs $\{\lambda,\lambda^c\}$ and keep, in the
circuit of Proposition~\ref{prop:promise}, only the terms for these pairs.
This circuit has $g$ composition gates. On a coloring in $\mathcal W$, it
is correct on both structures of the $2g$ tests whose colors lie in the
chosen pairs, and for every other color it answers zero at the read
edge, which is correct when the bit is one. Its success probability is
therefore at least $(1-\delta)(\frac12+g/2^k)\ge\frac12+g/2^k-\delta$.
\end{proof}

\subsection{Reflexive--transitive closure}\label{sec:closure}
\closure*
\begin{proof}
On the reference structure, \eqref{eq:table-off} gives $X;X=\one$ whenever
the off-diagonal support $S_X$ is nonempty. Hence $X^*=\one$ if
$S_X\ne\varnothing$ and $X^*=\id$ otherwise; in particular, closure
preserves color-determined relations, as condition~(i) of a test system requires. Let
$\lambda$ be a color, and put $J\coloneqq\Lambda\setminus\{\lambda\}$. If
$S_X=\varnothing$, then $\pred X\lambda=X$, and both
$(\pred X\lambda)^*$ and $\pred{(X^*)}\lambda$ equal $\id$. If $S_X$
contains a color in $J$, the table of \S\ref{sec:responses} gives
$\pred X\lambda;\pred X\lambda=\one$, so
$(\pred X\lambda)^*=\one=\pred{(X^*)}\lambda$. Thus a closure gate
notices the fold of $\lambda$ only when $S_X=\{\lambda\}$, so it notices
at most one fold. Theorem~\ref{thm:test-cost} and Lemma~\ref{lem:cor-cost}
give $2^k\le m_1(C)+2m_2(C)+g_*(C)\le2\gcost C+g_*(C)$ with the tests of
Lemma~\ref{lem:G-tests}.
\end{proof}
